\documentclass[10pt,a4paper]{amsart}
\usepackage{amsmath,amssymb,amsthm,graphicx}
\usepackage{verbatim}
\newenvironment{reference}{\comment}{\endcomment}
%\newenvironment{reference}{}{}
\newenvironment{slogan}{\comment}{\endcomment}
\newenvironment{history}{\comment}{\endcomment}

% For commutative diagrams we use Xy-pic
\usepackage[all]{xy}

% We use 2cell for 2-commutative diagrams.
\xyoption{2cell}
% Theorem environments.
%
\theoremstyle{plain}
\newtheorem{theorem}[subsection]{Theorem}
\newtheorem{proposition}[subsection]{Proposition}
\newtheorem{lemma}[subsection]{Lemma}

\theoremstyle{definition}
\newtheorem{definition}[subsection]{Definition}
\newtheorem{example}[subsection]{Example}
\newtheorem{exercise}[subsection]{Exercise}
\newtheorem{situation}[subsection]{Situation}

\theoremstyle{remark}
\newtheorem{remark}[subsection]{Remark}
\newtheorem{remarks}[subsection]{Remarks}

\numberwithin{equation}{subsection}

% Macros
%
\def\lim{\mathop{\mathrm{lim}}\nolimits}
\def\colim{\mathop{\mathrm{colim}}\nolimits}
\def\Spec{\mathop{\mathrm{Spec}}}
\def\Hom{\mathop{\mathrm{Hom}}\nolimits}
\def\Ext{\mathop{\mathrm{Ext}}\nolimits}
\def\SheafHom{\mathop{\mathcal{H}\!\mathit{om}}\nolimits}
\def\SheafExt{\mathop{\mathcal{E}\!\mathit{xt}}\nolimits}
\def\Sch{\mathit{Sch}}
\def\Mor{\mathop{\mathrm{Mor}}\nolimits}
\def\Ob{\mathop{\mathrm{Ob}}\nolimits}
\def\Sh{\mathop{\mathit{Sh}}\nolimits}
\def\NL{\mathop{N\!L}\nolimits}
\def\CH{\mathop{\mathrm{CH}}\nolimits}
\def\proetale{{pro\text{-}\acute{e}tale}}
\def\etale{{\acute{e}tale}}
\def\QCoh{\mathit{QCoh}}
\def\Ker{\mathop{\mathrm{Ker}}}
\def\Im{\mathop{\mathrm{Im}}}
\def\Coker{\mathop{\mathrm{Coker}}}
\def\Coim{\mathop{\mathrm{Coim}}}

% Boxtimes
%
\DeclareMathSymbol{\boxtimes}{\mathbin}{AMSa}{"02}

%
% Macros for moduli stacks/spaces
%
\def\QCohstack{\mathcal{QC}\!\mathit{oh}}
\def\Cohstack{\mathcal{C}\!\mathit{oh}}
\def\Spacesstack{\mathcal{S}\!\mathit{paces}}
\def\Quotfunctor{\mathrm{Quot}}
\def\Hilbfunctor{\mathrm{Hilb}}
\def\Curvesstack{\mathcal{C}\!\mathit{urves}}
\def\Polarizedstack{\mathcal{P}\!\mathit{olarized}}
\def\Complexesstack{\mathcal{C}\!\mathit{omplexes}}
% \Pic is the operator that assigns to X its picard group, usage \Pic(X)
% \Picardstack_{X/B} denotes the Picard stack of X over B
% \Picardfunctor_{X/B} denotes the Picard functor of X over B
\def\Pic{\mathop{\mathrm{Pic}}\nolimits}
\def\Picardstack{\mathcal{P}\!\mathit{ic}}
\def\Picardfunctor{\mathrm{Pic}}
\def\Deformationcategory{\mathcal{D}\!\mathit{ef}}
\usepackage{cleveref}
\crefname{originalreference}{original reference}{original references}
\Crefname{originalreference}{Original reference}{Original references}
\makeatletter
\newlabel{lemma-elementary-properties-change-of-field}{{07WA}{source}{}{sourcepointer-1}{}}
\newlabel{lemma-elementary-properties-change-of-field@cref}{{[originalreference][0][]07WA}{[1][0][]source}{}{}{}}
\makeatother
\crefname{originalreference}{original reference}{original references}
\Crefname{originalreference}{Original reference}{Original references}
\makeatletter
\newlabel{lemma-linear-functor}{{06I6}{source}{}{sourcepointer-2}{}}
\newlabel{lemma-linear-functor@cref}{{[originalreference][0][]06I6}{[1][0][]source}{}{}{}}
\makeatother
\crefname{originalreference}{original reference}{original references}
\Crefname{originalreference}{Original reference}{Original references}
\makeatletter
\newlabel{lemma-tangent-space-vector-space}{{06IH}{source}{}{sourcepointer-3}{}}
\newlabel{lemma-tangent-space-vector-space@cref}{{[originalreference][0][]06IH}{[1][0][]source}{}{}{}}
\makeatother
\crefname{originalreference}{original reference}{original references}
\Crefname{originalreference}{Original reference}{Original references}
\makeatletter
\newlabel{situation-change-of-fields}{{07W8}{source}{}{sourcepointer-4}{}}
\newlabel{situation-change-of-fields@cref}{{[originalreference][0][]07W8}{[1][0][]source}{}{}{}}
\makeatother
\title{The tangent space of the change of residue field of a predeformation category}
\author{The Stacks Project Authors}
\date{stacks-project:formal-defos.tex:local-snapshot-20261012; GNU FDL 1.2 or later}
\begin{document}
\maketitle
\noindent\textit{Editorial note.} All delivered mathematical text is the author's own, copied from the pinned Stacks Project chapter file formal-defos.tex; no statement or proof step was rewritten, reordered or removed.\par\medskip
\noindent\textit{Editorial note.} The counted proof is the author's complete argument: the reduction to the associated category in sets with the chain of equalities identifying the tangent space of the changed category, the generalization to trivial extensions by finite dimensional vector spaces, and the application of uniqueness of a linear structure.\par\medskip
\noindent\textit{Editorial note.} The label lemma-tangent-space-change-of-field names the counted statement, and the proof environment immediately following it in the source is delivered verbatim as its proof.\par\medskip
\noindent\textit{Editorial note.} The counted statement is carried by the author's own lemma environment in the chapter source.\par\medskip
\noindent\textit{Editorial note.} Pre-open gate: python operations/agent-authoring/tools/precard\_gate.py --source formal-defos.tex --label lemma-tangent-space-change-of-field --start 7598 --end 7632 --local-id 28029 returned RESULT=OK (L28029 free, occupancy FREE, public interval clear). Candidate taken from sources/stacks-formal-defos/candidates-halfA.json.\par\medskip
\noindent\textit{Editorial note.} The statement and proof are copied from the pinned chapter file; the cited inputs are kept as external pointers with their canonical Stacks tags.\par\medskip
\noindent\textit{Allowed standard prerequisites (editorial).}\begin{itemize}
\item Elementary properties of the change of residue field construction. Same chapter, the lemma collecting the elementary properties of the change of residue field of a category cofibred in groupoids. Supplies the reduction to the associated category in sets and the identification of the deformation situations.
\item Vector space structure on the tangent space and its description by the value at the trivial extension. Same chapter, the lemma providing the natural vector space structure on the tangent space and identifying its elements with the values of the functor on the trivial extension. Supplies the chain of equalities computing the tangent space of the changed category.
\item Uniqueness of a linear structure compatible with a functor. Same chapter, the lemma that a functor to the category of sets admits at most one structure of a linear functor. Supplies the conclusion that the identifications are compatible with the l-vector space structures.
\end{itemize}\medskip
\ifdefined\hypertarget\hypertarget{sourcepointer-1}{}\fi
\noindent\textit{Source reference.} Original reference 07WA: sources/stacks-formal-defos/pinned-original-source.tex, source lines 7539--7539. Prerequisite pointer invoked at proof line(s) 7611; it names a standard fact of the same chapter that the counted argument uses, and every step of the counted proof built on it is retained verbatim in the delivered body.\par\medskip
\ifdefined\hypertarget\hypertarget{sourcepointer-2}{}\fi
\noindent\textit{Source reference.} Original reference 06I6: sources/stacks-formal-defos/pinned-original-source.tex, source lines 3369--3369. Prerequisite pointer invoked at proof line(s) 7628; it names a standard fact of the same chapter that the counted argument uses, and every step of the counted proof built on it is retained verbatim in the delivered body.\par\medskip
\ifdefined\hypertarget\hypertarget{sourcepointer-3}{}\fi
\noindent\textit{Source reference.} Original reference 06IH: sources/stacks-formal-defos/pinned-original-source.tex, source lines 3669--3669. Prerequisite pointer invoked at proof line(s) 7615; it names a standard fact of the same chapter that the counted argument uses, and every step of the counted proof built on it is retained verbatim in the delivered body.\par\medskip
\ifdefined\hypertarget\hypertarget{sourcepointer-4}{}\fi
\noindent\textit{Source reference.} Original reference 07W8: sources/stacks-formal-defos/pinned-original-source.tex, source lines 7494--7494. Prerequisite pointer invoked at proof line(s) 7600; it names a standard fact of the same chapter that the counted argument uses, and every step of the counted proof built on it is retained verbatim in the delivered body.\par\medskip
\begin{lemma}
\label{lemma-tangent-space-change-of-field}
With notation and assumptions as in Situation \ref{situation-change-of-fields}.
Assume $\mathcal{F}$ is a predeformation category and
$\overline{\mathcal{F}}$ satisfies (S2). Then there is a
canonical $l$-vector space isomorphism
$$
T\mathcal{F} \otimes_k l \longrightarrow T\mathcal{F}_{l/k}
$$
of tangent spaces.
\end{lemma}

\begin{proof}
By Lemma \ref{lemma-elementary-properties-change-of-field} we may replace
$\mathcal{F}$ by $\overline{\mathcal{F}}$. Moreover we see that
$T\mathcal{F}$, resp.\ $T\mathcal{F}_{l/k}$ has a canonical
$k$-vector space structure, resp.\ $l$-vector space structure, see
Lemma \ref{lemma-tangent-space-vector-space}. Then
$$
T\mathcal{F}_{l/k} = \mathcal{F}_{l/k}(l[\epsilon])
= \mathcal{F}(k[l\epsilon]) = T\mathcal{F} \otimes_k l
$$
the last equality by Lemma \ref{lemma-tangent-space-vector-space}.
More generally, given a finite dimensional $l$-vector space $V$ we have
$$
\mathcal{F}_{l/k}(l[V]) = \mathcal{F}(k[V_k]) = T\mathcal{F} \otimes_k V_k
$$
where $V_k$ denotes $V$ seen as a $k$-vector space. We conclude that
the functors $V \mapsto \mathcal{F}_{l/k}(l[V])$ and
$V \mapsto T\mathcal{F} \otimes_k V_k$ are canonically identified
as functors to the category of sets. By Lemma \ref{lemma-linear-functor}
we see there is at most one way to turn either functor into an $l$-linear
functor. Hence the isomorphisms are compatible with the $l$-vector space
structures and we win.
\end{proof}

\end{document}
